\subsection{Arquitecturas de comunicación para robótica conectada}
\label{sec:arquitecturas_de_comunicacion_para_robotica_conectada}

El desarrollo del campo de la robótica móvil esta estrechamente vinculado al desarrollo de infraestructuras de comunicación capaces de soportar sistemas distribuidos y aplicaciones en tiempo real. Tradicionalmente, los robots móviles realizaban la mayor parte del procesamiento de forma local. Sin embargo, el creciente uso de sensórica, algoritmos de percepción basados en \gls{ia} y sistemas de control remoto ha incrementado significativamente las necesidades de computo.

En este contexto, las redes móviles \gls{5g} han emergido como una tecnología habilitadora de nuevas arquitecturas de robótica conectada. A diferencia de generaciones anteriores de redes móviles, el marco conceptual de \gls{5g} definido en IMT-2020 establece diferentes escenarios de servicio orientados a soportar aplicaciones con altos requisitos de capacidad, baja latencia y conectividad masiva de dispositivos \cite{3gppTS23501}. Estas capacidades permiten que robots móviles transmitan grandes volúmenes de datos procedentes de sensores y reciban comandos de control con tiempos de respuesta reducidos.

Desde el punto de vista arquitectónico, el sistema \gls{5g} se basa en un modelo denominado \gls{sba}, en el que las funciones de red se implementan como servicios modulares. Este enfoque facilita la virtualización de la red y la distribución flexible de funciones. Además, el sistema \gls{5g} introduce mecanismos como \textit{network slicing}, que permiten crear redes lógicas con características específicas para distintos servicios, y un modelo avanzado de \gls{qos} basado en identificadores \gls{5qi} para garantizar requisitos de latencia, fiabilidad y prioridad del tráfico. Asimismo, el estándar contempla el despliegue de \textit{Non-Public Networks} (NPN), que permiten operar infraestructuras \gls{5g} dedicadas para aplicaciones críticas \cite{3gppTS23501}.

Un ejemplo de red privadas \gls{5g} es la desplegada en el campus de la Universitat Politècnica de València, que proporciona cobertura interior y exterior mediante múltiples bandas de frecuencia y diferentes implementaciones del núcleo \gls{5g}, permitiendo la evaluación de diversos escenarios experimentales de conectividad \cite{ITEAM2025Private5G}. La red \gls{5gsa} del \gls{iteamupv} se describe con mayor detalle en la sección \ref{sec:capa_de_red_5g_privada}

El uso de redes \gls{5g} también facilita la transición hacia arquitecturas robóticas distribuidas, en las que determinadas funciones del sistema se ejecutan fuera de la plataforma robótica. Este enfoque permite reducir los requisitos de \textit{hardware} embarcado en el robot y mejorar la escalabilidad del sistema, aunque introduce nuevos retos relacionados con la latencia de comunicación, la estabilidad de los enlaces y la fiabilidad global de la red.

Diversos trabajos experimentales han explorado estas arquitecturas mediante plataformas de prueba que integran robots móviles con redes \gls{5g} privadas. En estos entornos, los robots transmiten flujos de datos sensoriales al \textit{edge} para su procesamiento y reciben información procesada para la toma de decisiones o la teleoperación. Los resultados muestran que el rendimiento del sistema depende de múltiples fuentes de latencia, incluyendo la transmisión de vídeo, el procesamiento de percepción asistido por \gls{ia} y el envío de comandos de control al robot \cite{Lozano2025DigitalTwin6G}.

En conjunto, estas investigaciones muestran que las redes \gls{5g} proporcionan la infraestructura necesaria para soportar arquitecturas robóticas distribuidas. No obstante, el comportamiento global del sistema depende de la interacción entre la red de comunicaciones, la arquitectura de procesamiento distribuido y el software robótico empleado. Por este motivo, el análisis del rendimiento de estos sistemas requiere considerar de forma conjunta las características de la red de comunicaciones y la distribución funcional de los componentes del sistema robótico.